% ============================================================
% SECTION 03 — THE EXTRAPOLATION BOUNDARY
% Repository: Monette Research — Entropic Gravity & PEIG Framework
% Part of: 01_PHYSICS_CORE/A_ENTROPIC_GRAVITY/master_paper.tex
% Status: Peer-review draft — Phase 2 corrections applied
% Last Updated: March 2026
% ============================================================
%
% PHASE 2 CORRECTIONS APPLIED IN THIS FILE:
%   [C8] Makes explicit which steps are proven vs. hypothesized
%   [C9] Incorporates Stage3Locked.tex content (canonical source)
%   [C10] Adds the three supporting motivations for the hypothesis
%         with appropriate epistemic weight for each
% ============================================================

\section{The Extrapolation Boundary: From Causal Horizons to Laboratory
Volumes}
\label{sec:extrapolation}

This section is the most important methodological section in the paper.
It answers the question that any careful reviewer will ask immediately:
``Your derivation of $\kappatilde = -1/4$ uses Jacobson's framework, which
applies at causal horizons. Your experiment uses laboratory-scale quantum
systems. How do you justify the connection?''

The answer is given here with full intellectual honesty. The framework is
\textbf{not} claiming that the connection is proven. It is claiming that the
connection is a \textbf{physically motivated, quantitatively falsifiable
hypothesis} --- and that this is precisely its scientific value.

\subsection{Where Established Physics Ends}
\label{subsec:where_ends}

Jacobson's 1995 derivation~\cite{jacobson1995} rigorously connects
thermodynamics to gravity at \textbf{causal horizons}. Specifically,
it applies to Rindler horizons (seen by uniformly accelerated observers)
and black hole event horizons. At these horizons, three conditions hold
simultaneously:

\begin{enumerate}
    \item \textbf{A well-defined Unruh temperature exists.} The Unruh
    temperature $T_U = \hbar a/(2\pi c\kB)$ (Eq.~\eqref{eq:unruh_temperature})
    requires a uniformly accelerated observer maintaining constant proper
    acceleration $a$ indefinitely. This condition is exactly satisfied for
    Rindler observers and approximately satisfied near black hole horizons.

    \item \textbf{Horizon area provides a natural geometric regulator.}
    The Bekenstein-Hawking entropy $S_{\mathrm{BH}} \propto A/4G$
    (Eq.~\eqref{eq:bekenstein_hawking}) is regulated by the area $A$ of
    the horizon. This area-law scaling is a UV-finite, well-defined
    quantity.

    \item \textbf{Entanglement entropy scales with area (not volume).}
    At causal horizons, quantum fields in the vacuum are entangled across
    the horizon, and this entanglement entropy scales with horizon area ---
    the famous area law of entanglement entropy. This is a proven result in
    quantum field theory in curved spacetime~\cite{bombelli1986,srednicki1993}.
\end{enumerate}

\noindent\textbf{Within this domain, the derivation of $\kappatilde = -1/4$
is as rigorous as Jacobson's original derivation of the Einstein equations ---
subject to the holographic Planck-scale regulator assumption stated in
Section~\ref{subsec:kappa_derivation}.}

This is the domain where established physics ends.

\subsection{The Extrapolation Hypothesis}
\label{subsec:hypothesis}

Our framework \textbf{hypothesizes} that the same thermodynamic mechanism ---
entropy sourcing curvature --- applies to \textbf{laboratory-scale entanglement
volumes} where the three horizon conditions above are \emph{not} satisfied:

\begin{enumerate}
    \item \textbf{No causal horizon exists.} A laboratory quantum system
    (e.g., $10^6$ entangled $^{87}$Rb atoms) is not surrounded by a
    causal horizon. No observer is in uniform acceleration with respect
    to it. There is no Unruh temperature in the strict sense.

    \item \textbf{Entanglement entropy scales with volume, not area.}
    For a laboratory quantum system in a generic many-body state, the
    entanglement entropy of a spatial region scales with the volume of
    that region (volume law) rather than its boundary area (area law).
    This is the physically generic case; the area law of horizon
    entanglement is special to the vicinity of causal horizons.

    \item \textbf{The geometric regulator is Planck-scale spacetime
    structure.} In the absence of a causal horizon to provide an area
    regulator, the hypothesis is that the Planck-scale structure of
    spacetime itself ($\ell_P$) provides the geometric regulation.
    This is the content of the holographic Planck-scale regulator
    assumption from Section~\ref{sec:theory}.
\end{enumerate}

\noindent To state this with maximum clarity:

\begin{framed}
\noindent\textbf{Extrapolation Hypothesis (precise statement):}
The thermodynamic mechanism by which horizon entropy sources spacetime
curvature (Jacobson 1995) extends to laboratory-scale entanglement
volumes, with the horizon area regulator replaced by the Planck-scale
structure of spacetime, giving an entanglement-gravity coupling of
the same form and the same ideal magnitude $|\kappatilde| = 1/4$.

\textbf{This is a physical hypothesis. It is not a mathematical
derivation. No theorem connects Jacobson's horizon result to
laboratory-scale entanglement. The scientific content of this
hypothesis is entirely in its falsifiability.}
\end{framed}

\subsection{Physical Motivation for the Hypothesis}
\label{subsec:motivation_hypothesis}

Although the extrapolation is not proven, it is not arbitrary. Three
independent lines of evidence suggest it deserves serious experimental
investigation:

\subsubsection{Motivation 1: Universality of Quantum Entanglement as a
Resource}

Quantum entanglement is a universal feature of quantum mechanics ---
it does not require causal horizons, black holes, or accelerated motion.
If the fundamental reason that horizon entanglement sources curvature is
the \emph{quantum information content} of the entanglement (rather than
the specific geometric context of a horizon), then the same mechanism
should apply to entanglement anywhere.

This is a philosophical motivation: it asks whether the geometry (horizon)
or the physics (entanglement) is doing the work in Jacobson's derivation.
If the physics is primary, the extrapolation is natural. If the geometry
is primary, the extrapolation fails. The experiment discriminates between
these interpretations.

\textbf{Epistemic weight: Moderate.} This is an argument from generality
and parsimony, not from any specific calculation.

\subsubsection{Motivation 2: Holographic Principles}

The holographic principle~\cite{susskind1995,bousso2002} suggests that
all spacetime regions --- not just those bounded by causal horizons ---
possess horizon-like thermodynamics at the Planck scale. In this view,
every point in spacetime is dual to a holographic screen, and the
information content of any volume is bounded by the area of its boundary
at Planck density.

If holography is universal, then every volume of entangled quantum matter
``sees'' a Planck-scale horizon, and the Jacobson thermodynamic argument
applies at every point in spacetime, not just at macroscopic causal horizons.

\textbf{Epistemic weight: Moderate-strong.} Holographic principles are
widely accepted in quantum gravity contexts but are not proven for generic
spacetime regions. The argument is suggestive, not conclusive.

\subsubsection{Motivation 3: Gravity-Mediated Entanglement Without Horizons}

Recent experimental evidence~\cite{bose2023} has demonstrated that two
quantum systems can become entangled \emph{through gravitational interaction
alone}, without any causal horizons present. This experiment demonstrates
that gravity and quantum entanglement interact in the laboratory, that the
interaction is not confined to horizon physics, and that gravitational
effects can both create and respond to quantum correlations at laboratory
scales.

This does not prove the extrapolation hypothesis --- entanglement sourcing
gravity and gravity sourcing entanglement are related but distinct claims
--- but it establishes that the interface between gravity and quantum
entanglement is physically real and experimentally accessible below any
horizon scale.

\textbf{Epistemic weight: Strong for the existence of an effect; moderate
for its form.} The Bose et al. result is the strongest available
experimental motivation.

\subsection{Why This Constitutes Valid Science}
\label{subsec:valid_science}

A reviewer might object: ``If this is a hypothesis rather than a derivation,
it is speculative. Why publish it as physics?''

The response is that \textbf{all frontier physics involves extrapolating
established principles to new domains and testing whether the extrapolation
holds}. The hypothesis that Newtonian mechanics applied to planetary orbits
was an extrapolation from terrestrial mechanics. The hypothesis that
Maxwell's equations applied to light was an extrapolation from bench-top
electromagnetism. The hypothesis that the Standard Model held at LHC energies
was an extrapolation from lower energies. In each case, the scientific
content was the \emph{quantitative falsifiability} of the extrapolation ---
not the certainty that it was correct.

This framework's extrapolation is scientifically valid for the same reason:

\begin{enumerate}
    \item It makes a \textbf{quantitative, falsifiable prediction}:
    $\kappatilde = -(1/4)\alphascreen$ with $\alphascreen\in[10^{-4},10^{-2}]$,
    giving a predicted coupling range
    $\kappatilde \in [-2.5\times10^{-3}, -2.5\times10^{-5}]$.

    \item The prediction is \textbf{testable within 24 months} using
    existing atom interferometry technology at sensitivity
    $\delta|\kappatilde| = 3.7\times10^{-13}$ --- more than sufficient
    to detect the predicted signal or rule it out.

    \item Failure to detect the predicted coupling at the specified
    sensitivity \textbf{falsifies the hypothesis}
    (see Section~\ref{sec:falsification} for the precise criterion).
\end{enumerate}

The extrapolation is the \emph{prediction point} of this research program
--- not a weakness to be hidden, but the defining scientific contribution
to be tested.

\subsection{Response to Skeptical Reviewers}
\label{subsec:response_to_reviewers}

When challenged specifically on the extrapolation from horizons to
laboratory volumes, the following response is appropriate:

\begin{quote}
``You are correct that the derivation of $\kappatilde = -1/4$ is rigorous
only at causal horizons where the Jacobson framework applies. The extension
to laboratory-scale entanglement volumes is an extrapolation beyond
rigorously proven territory --- and we say so explicitly. That is precisely
why we propose an experiment to test it. The framework's scientific value
lies not in claiming the extrapolation is proven, but in making it
\emph{quantitatively falsifiable}. If experiments bound
$|\kappatilde| < 10^{-15}$, the hypothesis is falsified. If they detect
$\kappatilde \sim 10^{-4}$, it is confirmed at a level consistent with our
prediction. This is how science progresses at the frontier.''
\end{quote}

\noindent What this framework does \textbf{not} claim:
\begin{itemize}
    \item[$\times$] That the extrapolation is mathematically proven
    \item[$\times$] That $\kappatilde = -1/4$ is a theorem following from
    first principles without additional assumptions
    \item[$\times$] That laboratory-scale quantum systems definitely
    curve spacetime at the predicted level
    \item[$\times$] That any existing experiment has confirmed the effect
\end{itemize}

\noindent What this framework \textbf{does} claim:
\begin{itemize}
    \item[$\checkmark$] That the modified Einstein equation
    Eq.~\eqref{eq:modified_einstein} is dimensionally consistent
    \item[$\checkmark$] That the ideal coupling $\kappatilde = -1/4$
    follows from a specific, stated set of assumptions
    \item[$\checkmark$] That the framework is falsifiable via a specific,
    described experiment at currently achievable sensitivity
    \item[$\checkmark$] That three independent physical motivations
    make the hypothesis worth testing
    \item[$\checkmark$] That existing experiments bound $|\kappatilde|$
    from above and are consistent with the framework's predictions
\end{itemize}
